\subsection{章末確認問題}
\begin{enumerate}[leftmargin=2em]
\item 欠陥と障害の違いを説明せよ。なぜテストは欠陥ではなく障害を直接観測することが多いのか。
\item テストケースに入力値以外の情報が必要な理由を説明せよ。
\item テストオラクルとは何か。オラクル問題が起きやすいシステムの例を一つ挙げよ。
\item 単体テスト、結合テスト、システムテスト、受入テストを、対象と目的の観点から比較せよ。
\item 回帰テストがアジャイルや DevOps で重要になる理由を説明せよ。
\item 同値分割と境界値分析の違いを、年齢入力の例で説明せよ。
\item 仕様ベース技法と構造ベース技法はなぜ補完関係にあるのか。
\item 探索的テストが有効になる状況を説明せよ。
\item ミューテーションスコアは何を表すか。値が高いと何が言えるか。
\item テスト完了をカバレッジだけで判断してはいけない理由を説明せよ。
\item シフトレフトの考え方を、従来型プロセスと比較して説明せよ。
\item 「新興技術をテストする」と「新興技術を使ってテストする」の違いを説明せよ。
\end{enumerate}

\begin{pointbox}{解答の方向性}
1 は、原因としての欠陥と観測される結果としての障害を区別する。2 は、状態、環境、手順、期待結果が振る舞いに影響することを述べる。3 は、結果判定の仕組みであり、機械学習や推薦などで難しい。4 は対象の大きさと利用者視点への近さで整理する。5 は変更が頻繁で、既存機能の破壊を早く見つける必要があることを述べる。6 以降は、本文の「選択基準」「十分性」「カバレッジ」「リスク」「費用」「再現性」「目的に応じた技法選択」という語を使って説明できる。
\end{pointbox}
